$\displaystyle \sum_{i=1}^{n} \frac{1}{16i^2-1}$

\nopagebreak
\begin{align*}
\text{Factorizamos el denominador: } 16i^2 - 1 &= (4i-1)(4i+1) \\
\text{Fracciones parciales: } \frac{1}{(4i-1)(4i+1)} &= \frac{A}{4i-1} + \frac{B}{4i+1} \\
1 &= A(4i+1) + B(4i-1) \\
\text{Si } i=1/4 \implies 1 &= 2A \implies A=1/2 \\
\text{Si } i=-1/4 \implies 1 &= -2B \implies B=-1/2 \\
\text{La suma es: } \sum_{i=1}^{n} \frac{1}{2} \left( \frac{1}{4i-1} - \frac{1}{4i+1} \right) &= \frac{1}{2} \left[ \left( \frac{1}{3} - \frac{1}{5} \right) + \left( \frac{1}{5} - \frac{1}{9} \right) + \dots + \left( \frac{1}{4n-1} - \frac{1}{4n+1} \right) \right] \\
\text{Es telescópica, se cancelan términos intermedios:} \\
&= \frac{1}{2} \left( \frac{1}{3} - \frac{1}{4n+1} \right)
\end{align*}

\[
\boxed{\displaystyle \frac{1}{2} \left( \frac{1}{3} - \frac{1}{4n+1} \right) = \frac{2n-1}{3(4n+1)}}
\]
